\section{Introduction}
\label{sec:intro}

The deployment of deep neural networks onto edge devices, such as mobile phones, smart sensors, and automotive processors, is strictly constrained by computational, memory, and energy budgets. At the same time, real-world edge applications frequently require the execution of multiple specialized tasks. Running a separate model for each task is highly inefficient and often impossible under strict on-device memory constraints. To mitigate this, weight-space model merging, such as \emph{Task Arithmetic}~\cite{ilharco22} and \emph{Model Soups}~\cite{wortsman22}, has emerged as a powerful paradigm. By linearly combining the weights of multiple specialized ``task experts'' fine-tuned from a shared pre-trained backbone, a single multi-task model can be constructed. This approach incurs absolutely zero computational or memory overhead at inference time compared to running a single model.

To meet the severe hardware limits of edge accelerators, these merged models must undergo aggressive post-training quantization (PTQ) to low-bit integers (e.g., 4-bit or 8-bit representation)~\cite{jacob18quant, gholami21survey}. Traditional full-precision merging techniques like \emph{AdaMerging}~\cite{yang23adamerging} do not account for quantization during optimization, resulting in severe accuracy drops when the model is subsequently quantized. To address this, quantization-aware model merging frameworks such as \emph{Q-Merge}~\cite{qmerge} use the Straight-Through Estimator (STE)~\cite{bengio13ste} to optimize layer-wise merging coefficients directly under a simulated quantization operator.

\begin{figure}[t]
\centering
\includegraphics[width=\linewidth]{fig1.png}
\caption{Cross-schema fidelity accuracy retention (\%) of different model merging methods under robust 8-bit quantization across 5 heterogeneous target hardware operators. Standard Q-Merge optimized strictly under a Symmetric Per-Channel operator suffers from notable performance degradation when deployed onto mismatched schemas (e.g., dropping to 45.90\% under Symmetric Per-Tensor). Our proposed \textbf{OmniMerge} stochastically co-optimizes across multiple operators, maintaining robust and high performance across all target environments.}
\label{fig:cross_schema_intro}
\end{figure}

However, as a pragmatic investigation of real-world edge deployment reveals, optimizing strictly under a single simulated operator is a dangerous simplification. In practice, edge-computing fleets are highly heterogeneous. Different hardware ASICs, digital signal processors (DSPs), and runtime compilers (such as Google Edge TPU, Apple Neural Engine, or NVIDIA TensorRT) implement distinct, incompatible post-training quantization standards. These standards range from symmetric per-channel to asymmetric per-tensor quantization, each utilizing different rounding boundaries, scale factors, and zero-point alignments.

In this work, we demonstrate that optimizing merging coefficients strictly under a single quantization operator (as done by Q-Merge) induces a phenomenon we term \textbf{Cross-Schema Performance Degradation}. Because first-order optimization via STE adjusts coefficients to fit the exact, discrete rounding grid of the training operator, the model overfits to those specific discretization boundaries. When the resulting merged model is deployed on a mismatched target compiler or hardware ASIC, this boundary-overfitting causes the gradient-derived coefficients to fail, leading to significant performance drops relative to optimal multi-schema ensembling.

To resolve this critical deployment bottleneck, we propose \textbf{OmniMerge}, a training-free, multi-schema stochastic co-optimization framework that produces robust, hardware-invariant merging coefficients. Guided by a pragmatic philosophy, OmniMerge requires \textbf{zero hardware metadata}, adds \textbf{no extra test-time compute}, and introduces \textbf{zero inference-time latency or memory overhead} because the optimized coefficients are used to compile a standard quantized model.

OmniMerge achieves this robust generalization via two core techniques:
\begin{itemize}
    \item \textbf{Stochastic Operator Sampling (SOS):} During test-time adaptation, rather than utilizing a single static quantization operator, the active quantization operator is stochastically sampled from a discrete pool of hardware-relevant schemas at each optimization step. This acts as parameter-space data augmentation, preventing the continuous coefficients from overfitting to any single rounding grid.
    \item \textbf{Scale and Zero-Point Noise Perturbation (SZNP):} To enhance generalization to unseen runtime compilers, we perturb the scale factors and zero-point offsets dynamically with multiplicative and additive Gaussian noise. This smooths out the local loss landscape, helping the optimizer avoid brittle local minima.
\end{itemize}

We evaluate OmniMerge using a Vision Transformer backbone (\texttt{ViT-Tiny})~\cite{dosovitskiy20vit} merged across four diverse classification tasks (MNIST, FashionMNIST, CIFAR-10, and SVHN) under robust 8-bit quantization. As shown in Figure~\ref{fig:cross_schema_intro}, OmniMerge successfully closes the cross-schema generalization gap, achieving up to \textbf{50.78\%} average accuracy across five mismatched target operators. This represents a significant breakthrough for practical edge-AI deployment, ensuring that a single optimization pass can produce a merged checkpoint capable of running on any heterogeneous hardware target.
